\section{Parameter estimates and a comparison trajectory}
\label{sec:parameter-trajectory}

The correction in Section~\ref{sec:projected-correction} was obtained
at each fixed \(a\). We now control its parameter derivative. The
derivative estimates still allow repeated line factors. The scalar
trajectory constructed at the end of this section requires
\(L\not\simeq M\).

\subsection{Differentiation in a fixed unitary bundle}

Write \(H_a=H_0T_a\), with \(T_a\) as in
\eqref{eq:compact-metric}. The isometry
\(U_a=T_a^{1/2}:(V,H_a)\to(V,H_0)\) satisfies
\(a\partial_aU_a=\frac12U_ak_a\), since \(T_a\) and its
parameter derivative commute. For a moving Hermitian section
define its derivative in this fixed unitary identification by
\begin{equation}\label{eq:unitary-parameter-derivative}
 \dot s=a\partial_as+\tfrac12[k_a,s].
\end{equation}
A dot on an operator has the same meaning. The base volume and
the Hilbert inner product in the identified bundle are fixed.

\begin{lemma}\label{lem:background-parameter}
The explicit family satisfies
\begin{equation}\label{eq:background-parameter-bounds}
 \|k_a\|_{X_a}\le C,\qquad
 \|\dot\Phi_a\|_{Y_a}\le Ca,\qquad
 \|\dot{\mathcal L}_a\|_{X_a\to Y_a}\le C.
\end{equation}
Moreover
\begin{equation}\label{eq:parameter-linearization}
 \dot\Phi_a=\mathcal L_ak_a,\qquad
 \|Q_ak_a\|_{X_a}\le Ca\ell(a).
\end{equation}
\end{lemma}

\begin{proof}
On the annuli the identification by \(U_a\) fixes the unitary
graded frame: \(U_aF_a=F_0\), because its two multipliers are
\[
 w^{1/2}u_a^{-1/2}=u_0^{-1/2},\qquad
 w^{-1/2}u_a^{1/2}=u_0^{1/2}.
\]
Thus ordinary parameter differentiation of the normalized
graded matrices is exactly \eqref{eq:unitary-parameter-derivative}.
Applying \(a\partial_a\) to the explicit rational factors
with denominator \(a+u_0\) preserves their symbol estimates.
In particular, the cancelled residual formulas in
Proposition~\ref{prop:residual}, including their spatial
derivatives, remain \(O(a r^{-2-j})\).
The expression \eqref{eq:parameter-direction} similarly
has uniformly bounded scaled derivatives. This argument
differentiates the formulas and their cancellations, not
an unspecified remainder estimate.

On the core use \(z=\sqrt a\,y\). The base and line data
have vanishing first jets. Their integral second-order
Taylor remainders, and every fixed number of \(y\) derivatives,
are \(O(a)\), also after a logarithmic derivative at fixed
\(y\). Differentiation at fixed \(z\) is
\(a\partial_a-\frac12y\cdot\partial_y\) in these expressions.
Both terms preserve the bounds on a fixed enlarged ball.
The flat leading contracted curvature is identically zero.
Its physical residual and its logarithmic parameter derivative
are therefore bounded in every fixed \(y\) norm. The unitary
normalization and its derivative are uniformly bounded there.
This proves the first two estimates in
\eqref{eq:background-parameter-bounds} on the core as well.
The fixed exterior follows directly from smooth formulas.

The unitary variation of the Chern connection is the one-form
\[
 B_a=\tfrac12(D'_ak_a-D''_ak_a),
\]
acting by commutator on endomorphisms. It follows by
differentiating the Chern change formula and \(U_a\).
The bounds just proved give
\(|D_a^jB_a|\le C_j\rho_a^{-1-j}\).
Differentiation of \(\frac12D_a^*D_a\) has first-order
terms with coefficient \(B_a\) and zeroth-order terms
with one derivative of \(B_a\). The scaled Hölder
coefficient bounds prove the third estimate in
\eqref{eq:background-parameter-bounds}.
There is no derivative of the inner product in this
unitary identification. Finally the centered linearization
in the preceding section, applied to the metric velocity
\(k_a\), gives \(\dot\Phi_a=\mathcal L_ak_a\).
Apply \(G_a\) and \eqref{eq:log-inverse} to obtain
the second assertion in \eqref{eq:parameter-linearization}.
\end{proof}

\begin{proposition}\label{prop:correction-parameter}
The family \(v_a\) is smooth for \(a>0\), and
\begin{equation}\label{eq:correction-parameter}
 \|\dot v_a\|_{X_a}\le Ca\ell(a)^3.
\end{equation}
If the velocity of \(\widehat H_a\) is centered to \(H_a\) by
\[
 \widehat k_a^{\,c}
 =e^{v_a/2}\bigl(\widehat H_a^{-1}a\partial_a\widehat H_a\bigr)
     e^{-v_a/2},
\]
then
\begin{equation}\label{eq:corrected-tangent}
 \|\widehat k_a^{\,c}-k_a\|_\infty\le Ca\ell(a)^3.
\end{equation}
\end{proposition}

\begin{proof}
Use the fixed resolvent circle \(|z|=c/2\). Its distance from
both the low spectrum and the complementary spectrum is
bounded below. On that circle,
\[
 \|(z-\mathcal L_a)^{-1}\|_{Y_a\to X_a}\le C\ell(a).
\]
For the low summands this follows from
\eqref{eq:low-projection}. On the complement, multiply the
heat representation by \(e^{zt}\). Its short-time bound is
unchanged, and its large-time tail still decays exponentially
because \(\operatorname{Re}z\le c/2\). The kernel and scaled
Schauder arguments in Lemma~\ref{lem:weighted-inverse} apply.
The resolvent derivative and
\eqref{eq:background-parameter-bounds} consequently give
\begin{equation}\label{eq:projection-derivative}
 \|\dot\Pi_a\|_{Y_a\to X_a}\le C\ell(a)^2.
\end{equation}
Both resolvents map \(Y_a\) to \(X_a\); the middle
operator derivative maps \(X_a\) to \(Y_a\).
The bounded embedding \(X_a\to Y_a\) permits use of
the same estimate on \(X_a\). Differentiating
\(\Pi_a^2=\Pi_a\) shows that \(\dot\Pi_a\) exchanges
the low and complementary spaces; in particular
\(\Pi_a\dot\Pi_a\Pi_a=0\). No separation of eigenvalues
within the low cluster is required.

With \(v\) held fixed in the unitary identification, the
parameter derivative of the quadratic remainder obeys
\begin{equation}\label{eq:nonlinear-parameter}
 \|\dot N_a^{\,\mathrm{par}}(v)\|_{Y_a}\le C\|v\|_{X_a}^2.
\end{equation}
To see this, differentiate the finite curvature formula
after subtracting its value and first Fréchet derivative.
The scaled coefficients and their required parameter
derivatives are uniformly bounded by the explicit annular
and core calculations above. The second Fréchet derivative,
and its parameter derivative, are therefore uniformly
bounded on a fixed small ball. Matrix Taylor's formula
proves \eqref{eq:nonlinear-parameter}. The estimate includes
second-jet terms. The logarithmic derivative of \(\rho_a\)
is bounded, so the statement is in the same scaled norms
at the parameter under consideration.

The complementary derivative of the equation is invertible:
\begin{equation}\label{eq:perturbed-inverse}
 \left\|\bigl((\mathcal L_a+Q_aDN_a(v_a))|_{Q_aX_a}\bigr)^{-1}
       \right\|_{Q_aY_a\to Q_aX_a}\le C\ell(a).
\end{equation}
Indeed the perturbation times \(G_a\) is \(O(a\ell(a)^2)\),
so a Neumann series applies. On compact positive parameter
intervals, the projection and coefficients are smooth.
A local smooth identification of \(Q_a\) makes the
ordinary implicit function theorem applicable. Uniqueness
in Theorem~\ref{thm:projected-correction} identifies the local
solutions and proves smoothness for \(a>0\).

Put \(\mathcal F_a(v)=\Phi_a+N_a(v)\).
Differentiating \(\Pi_av_a=0\) gives
\[
 \Pi_a\dot v_a=-\dot\Pi_av_a,
 \qquad \|\Pi_a\dot v_a\|_{X_a}\le Ca\ell(a)^3.
\]
Differentiating
\(\mathcal L_av_a+Q_a\mathcal F_a(v_a)=0\) and projecting
onto \(Q_a\) gives
\begin{align*}
 (\mathcal L_a+Q_aDN_a(v_a))Q_a\dot v_a
 ={}&-Q_a\dot{\mathcal L}_av_a
      +Q_a\dot\Pi_a\mathcal F_a(v_a)-Q_a\dot\Phi_a\\
    &-Q_a\dot N_a^{\,\mathrm{par}}(v_a)
      -Q_aDN_a(v_a)\Pi_a\dot v_a.
\end{align*}
The five terms have \(Y_a\) norms
\[
 O(a\ell(a)),\quad O(a\ell(a)^2),\quad O(a),\quad
 O(a^2\ell(a)^2),\quad O(a^2\ell(a)^4),
\]
respectively. Their sum is \(O(a\ell(a)^2)\).
Applying \eqref{eq:perturbed-inverse} proves
\eqref{eq:correction-parameter}.

For the tangent, differentiation of the exponential yields
the exact identity
\begin{equation}\label{eq:centered-tangent-formula}
 \widehat k_a^{\,c}
 =\cosh(\operatorname{ad}_{v_a}/2)k_a+
   \frac{\sinh(\operatorname{ad}_{v_a}/2)}
        {\operatorname{ad}_{v_a}/2}\,\dot v_a.
\end{equation}
Substitute
\(a\partial_av_a=\dot v_a-\frac12[k_a,v_a]\)
in the integral derivative formula for \(e^{v_a}\);
the two \(k_a\) terms combine to the displayed hyperbolic
cosine. The quotient is defined by its entire power
series at zero. Its operator norm is bounded on the
small supremum ball, and the cosine differs from the
identity by \(O(|v_a|^2)\).
Estimates~\eqref{eq:corrected-bounds} and
\eqref{eq:correction-parameter} prove
\eqref{eq:corrected-tangent}.
\end{proof}

\subsection{The single small eigenvalue}

From now on assume \(L\not\simeq M\). Set
\begin{equation}\label{eq:scalar-coefficient}
 e_a=\Pi_ak_a,\qquad b_0=\frac{2\pi^2\sigma}{v_0}.
\end{equation}
The section \(e_a\) is nonzero for small \(a\): its mass tends
to \(2v_0\) by \eqref{eq:parameter-pairing} and
\eqref{eq:parameter-linearization}.

\begin{proposition}\label{prop:sharp-small-eigenvalue}
There is a smooth positive scalar \(b(a)\), for \(a\)
sufficiently small, such that
\begin{equation}\label{eq:remaining-scalar-curvature}
 \mathcal M_a(v_a)=b(a)e_a,\qquad b(a)=b_0a+o(a).
\end{equation}
For the real connection operators of \(H_a\) and
\(\widehat H_a\),
\begin{equation}\label{eq:sharp-small-eigenvalue}
 \lambda_1(\mathcal L_a)=b_0a+o(a),\qquad
 \lambda_1(\widehat{\mathcal L}_a)=b_0a+o(a),\qquad
 \lambda_2(\widehat{\mathcal L}_a)\ge c'>0.
\end{equation}
\end{proposition}

\begin{proof}
The remaining centered curvature belongs to the
one-dimensional space spanned by \(e_a\), defining \(b(a)\).
Pairing with \(k_a\) and using self-adjointness gives
\begin{align*}
 b(a)\|e_a\|_2^2
 &=\langle\Phi_a,k_a\rangle
   +\langle v_a,\mathcal L_ak_a\rangle
   +\langle N_a(v_a),k_a\rangle\\
 &=4\pi^2\sigma a+o(a)+O(a^2\ell(a))+O(a^2\ell(a)^2).
\end{align*}
Here \(\|\mathcal L_ak_a\|_1=\|\dot\Phi_a\|_1=O(a)\).
Since \(\|e_a\|_2^2\to2v_0\), this proves
\eqref{eq:remaining-scalar-curvature}, including positivity.

We compute the eigenvalue separately. Put \(J=2P-I\) and
\(\zeta=u_0/(u_0+a)\), so \(k_a=\zeta J\) and \(u_a=a\zeta\).
In the unitary splitting, the diagonal derivative of \(\zeta\)
and off-diagonal derivative of \(J\) are orthogonal. Hence
\begin{equation}\label{eq:parameter-energy-density}
 |D_a''k_a|^2
 =\frac{2a^2|\partial u_0|^2+4a^2u_0^4|\beta|_0^2}{(u_0+a)^4}
 =\frac{4a^2u_0e_0-2a^2|\partial u_0|^2}{(u_0+a)^4}.
\end{equation}
The second equality uses \eqref{eq:full-inclusion-energy}
for \(H_0\). At a normalized reduced point,
\[
 u_0=\lambda_pr^2+O(r^4),\qquad
 e_0(p)=4\lambda_p,\qquad
 |\partial u_0|^2=2\lambda_p^2r^2+O(r^4).
\]
The density in \eqref{eq:parameter-energy-density} is
nonnegative and bounded by \(Ca^2r^2/(a+r^2)^4\).
After \(z=\sqrt a\,y\) and division of its integral by \(a\),
the dominating function is
\(C|y|^2/(1+|y|^2)^4\), integrable in four real dimensions.
The leading integral at \(p\) is
\[
 12\pi^2a^2\lambda_p^2
 \int_0^\infty\frac{s^2}{(a+\lambda_ps)^4}\,\dd s
 =\frac{4\pi^2a}{\lambda_p},
\]
since \(\int_0^\infty t^2(1+t)^{-4}\,\dd t=1/3\).
The fixed exterior is \(O(a^2)\). Thus
\(\mathfrak q_a(k_a,k_a)=4\pi^2\sigma a+o(a)\).
The complementary energy is
\[
 \mathfrak q_a(Q_ak_a,Q_ak_a)
 =\langle Q_ak_a,\mathcal L_ak_a\rangle
 =O(a^2\ell(a))=o(a).
\]
As \(e_a\) spans the first eigenspace and
\(\|e_a\|_2^2\to2v_0\), this proves the reference eigenvalue
equivalent.

Under the isometry \(e^{v_a/2}\), the corrected unitary
connection differs from \(D_a\) by a one-form \(A_a\) with
\(|A_a|\le C|D_av_a|\). The weighted derivative bound gives
\[
 \|A_a\|_4\le Ca\ell(a)^{5/4},\qquad
 \int_X\rho_a^{-4}\,\dd V\le C\ell(a).
\]
For every real Hermitian trace-free section \(s\), the
uniform Sobolev inequality yields
\[
 \|A_as\|_2
 \le Ca\ell(a)^{5/4}
       \bigl(\mathfrak q_a(s,s)+\|s\|_2^2\bigr)^{1/2}.
\]
Expansion of the corrected connection energy and the
product inequality with parameter \(\sqrt a\) imply
\begin{equation}\label{eq:corrected-form-comparison}
 |\widehat{\mathfrak q}_a(s,s)-\mathfrak q_a(s,s)|
 \le C\sqrt a\,\mathfrak q_a(s,s)
      +Ca^{3/2}\ell(a)^{5/2}\|s\|_2^2.
\end{equation}
A term of order \(a^{3/2}\ell(a)^{5/2}\mathfrak q_a\)
is absorbed into the first term for small \(a\).
The mass forms are identical under the isometry.
Min--max, the reference equivalent and its complementary
gap prove the remaining assertions in
\eqref{eq:sharp-small-eigenvalue}.
\end{proof}

\subsection{A trajectory with integrable flow error}

Choose a sufficiently small fixed \(a_{\mathrm{in}}>0\) and solve
\begin{equation}\label{eq:comparison-ode}
 a'(T)=-2a(T)b(a(T)),\qquad a(0)=a_{\mathrm{in}}.
\end{equation}
Define \(\overline H(T)=\widehat H_{a(T)}\). This is a curve
of smooth metrics on the original compact bundle.

\begin{proposition}\label{prop:comparison-trajectory}
The curve exists for all \(T\ge0\), and
\begin{equation}\label{eq:comparison-asymptotic}
 a(T)\sim\frac1{2b_0T},\qquad
 \overline Z(T)\sim\frac{v_0^2}{4\pi^2\sigma T}.
\end{equation}
Its metric-flow error
\[
 \mathcal R(T)=\overline H^{-1}\partial_T\overline H
                 +2\Phi_{\overline H}
\]
obeys
\begin{equation}\label{eq:integrable-flow-error}
 \|\mathcal R(T)\|_{\infty,\overline H}
 \le Ca(T)^2\ell(a(T))^3,\qquad
 \int_T^\infty\|\mathcal R(t)\|_\infty\,\dd t
 \le Ca(T)\ell(a(T))^3.
\end{equation}
\end{proposition}

\begin{proof}
For small \(a\), \(b(a)/a\) lies between two positive
constants. The solution stays positive, decreases to zero
and exists for all nonnegative time. Its reciprocal
satisfies
\[
 (1/a)'=2b(a)/a\longrightarrow2b_0.
\]
Integration proves the first equivalent in
\eqref{eq:comparison-asymptotic}; the corrected mass
in \eqref{eq:corrected-energy} proves the second.
Centering the error gives the exact expression
\[
 e^{v_a/2}\mathcal R e^{-v_a/2}
     =2b(a)\bigl(e_a-\widehat k_a^{\,c}\bigr).
\]
The term in parentheses has supremum norm
\(O(a\ell(a)^3)\) by
\eqref{eq:parameter-linearization} and
\eqref{eq:corrected-tangent}. Since \(b(a)=O(a)\),
this proves the first estimate in
\eqref{eq:integrable-flow-error}. For the second, substitute
\(\dd T=-\dd a/(2ab(a))\) and use \(b(a)\ge ca\).
The tail is bounded by
\[
 C\int_0^{a(T)}(1+|\log x|)^3\,\dd x
 \le Ca(T)\ell(a(T))^3.
\]
\end{proof}

The coefficient and time scale in this proposition belong so far
to the constructed curve. The next section proves convergence of
every actual flow with the fixed determinant to that curve. In
particular, integrability of the error alone will first give
bounded metric distance; the small positive eigenvalue supplies
the additional decay.
